\paragraph{A bound for the Chilean cohorts}
A designer who knows the April enrollments can add at least 642 historical hits while keeping every baseline recipient seated on all 1,024 draws per cohort. That is 328 more than the 314 possible when everyone must keep the same course and seat type. We checked each constructed set of menus by re-clearing the national market on every draw. At most 4,620 extra hits are possible. The upper bound rules out options at which someone must be rejected or must displace a pupil who has nowhere else to go. The unrestricted switch gained 3,590 hits. Whether a seat-retaining policy can match that gain remains open.

\begin{table}[ht]\centering\small
\caption{Additional pupils whose April school can appear on their menu. Both endpoints refer to the same seat-retention problem. The first and last numeric columns give exact optima under the stronger contract requirement and with no assignment constraint.}\label{tab:retention-bounds}
\begin{tabular}{@{}lrrrr@{}}\toprule
 & Same contract & \multicolumn{2}{c}{Keep a seat} & Unconstrained\\
\cmidrule(lr){3-4}
Application year & & Achieved & Upper bound & \\\midrule
2022 & 96 & 170 & 1,350 & 1,425\\
2023 & 111 & 199 & 1,564 & 1,633\\
2024 & 107 & 273 & 1,706 & 1,765\\
\midrule
Total & 314 & 642 & 4,620 & 4,823\\\bottomrule
\end{tabular}
\end{table}
